% GENERATED FILE. Edit canonical lessons, then run npm run generate:books.
% canonical-source-hash: fnv1a64:eda85d4a8df532b3
% canonical-lessons: UR-C76-minat, UR-C76-aj, UR-C76-kal, UR-C76-abhi, UR-C76-phir

\chapter{Minute, Today, Yesterday, Right now, Then}
\label{ch:ur-minute-today-yesterday-right-now-then}
\hlchaptermodality{76}

\begin{chapteropening}
\textbf{By the end of this chapter:} \emph{I can say a minute, today, yesterday, right now and then, again.}

Complete the last lesson of chapter 76: I can say a minute, today, yesterday, right now and then, again.
\end{chapteropening}

\section[\texorpdfstring{\ur{منٹ} (minaṭ)}{منٹ (minaṭ)}]{\ur{منٹ} (minaṭ) — a minute}

\label{lesson:UR-C76-minat}



\begin{quote}
Before the new one: say the Urdu for a year, then the Urdu for an hour.
\end{quote}

\subsection*{You'll want to know: \ur{منٹ}}
\textbf{\ur{منٹ}} — \emph{minaṭ} — \textquotedblleft{}a minute\textquotedblright{}.

\begin{cousinweb}
\textbf{\ur{منٹ}} (\emph{minaṭ}) is a loan from English \emph{minute}.
\end{cousinweb}

\begin{grammarlens}[title={What you've built}]
One more word for this chapter.
\end{grammarlens}

\subsection*{Your turn}
\begin{itemize}
\raggedright
  \item \emph{Say it:} \emph{minaṭ}
  \item \emph{Say it:} \emph{minaṭ}, once more
  \item \emph{Say it:} say \emph{minaṭ}
  \item \emph{Recall:} say the Urdu for a year, then the Urdu for an hour, then say \emph{minaṭ} again
\end{itemize}

\begin{tcolorbox}[breakable,colback=teal!4,colframe=teal!35!black,title={Before you move on}]
What does \ur{منٹ} mean? (\textquotedblleft{}A minute\textquotedblright{}.) Say it once more.
\end{tcolorbox}
\section[\texorpdfstring{\ur{آج} (āj)}{آج (āj)}]{\ur{آج} (āj) — today}

\label{lesson:UR-C76-aj}



\begin{quote}
Before the new one: say the Urdu for an hour, then the Urdu for a minute.
\end{quote}

\subsection*{You'll want to know: \ur{آج}}
\textbf{\ur{آج}} — \emph{āj} — \textquotedblleft{}today\textquotedblright{}.

\begin{cousinweb}
\textbf{\ur{آج}} (\emph{āj}) is a tadbhava, a word inherited from Sanskrit and reshaped by sound change. It comes from \emph{adya}, \textquotedblleft{}today\textquotedblright{} (CDIAL 242).
\end{cousinweb}

\begin{grammarlens}[title={What you've built}]
One more word for this chapter.
\end{grammarlens}

\subsection*{Your turn}
\begin{itemize}
\raggedright
  \item \emph{Say it:} \emph{āj}
  \item \emph{Say it:} \emph{āj}, once more
  \item \emph{Say it:} say \emph{āj}
  \item \emph{Recall:} say the Urdu for an hour, then the Urdu for a minute, then say \emph{āj} again
\end{itemize}

\begin{tcolorbox}[breakable,colback=teal!4,colframe=teal!35!black,title={Before you move on}]
What does \ur{آج} mean? (\textquotedblleft{}Today\textquotedblright{}.) Say it once more.
\end{tcolorbox}
\section[\texorpdfstring{\ur{کل} (kal)}{کل (kal)}]{\ur{کل} (kal) — yesterday; tomorrow}

\label{lesson:UR-C76-kal}



\begin{quote}
Before the new one: say the Urdu for a minute, then the Urdu for today.
\end{quote}

\subsection*{You'll want to know: \ur{کل}}
\textbf{\ur{کل}} — \emph{kal} — \textquotedblleft{}yesterday; tomorrow\textquotedblright{}.

\begin{cousinweb}
\textbf{\ur{کل}} (\emph{kal}) is a tadbhava, a word inherited from Sanskrit and reshaped by sound change. It comes from \emph{kalya}, \textquotedblleft{}dawn; at dawn, tomorrow\textquotedblright{} (CDIAL 3104). It is built on \emph{kāla}, \textquotedblleft{}time\textquotedblright{}; one word serves for both \textquotedblleft{}yesterday\textquotedblright{} and \textquotedblleft{}tomorrow\textquotedblright{}.
\end{cousinweb}

\begin{grammarlens}[title={What you've built}]
One more word for this chapter.
\end{grammarlens}

\subsection*{Your turn}
\begin{itemize}
\raggedright
  \item \emph{Say it:} \emph{kal}
  \item \emph{Say it:} \emph{kal}, once more
  \item \emph{Say it:} say \emph{kal}
  \item \emph{Recall:} say the Urdu for a minute, then the Urdu for today, then say \emph{kal} again
\end{itemize}

\begin{tcolorbox}[breakable,colback=teal!4,colframe=teal!35!black,title={Before you move on}]
What does \ur{کل} mean? (\textquotedblleft{}Yesterday; tomorrow\textquotedblright{}.) Say it once more.
\end{tcolorbox}
\section[\texorpdfstring{\ur{ابھی} (abhī)}{ابھی (abhī)}]{\ur{ابھی} (abhī) — right now}

\label{lesson:UR-C76-abhi}



\begin{quote}
Before the new one: say the Urdu for today, then the Urdu for yesterday; tomorrow.
\end{quote}

\subsection*{You'll want to know: \ur{ابھی}}
\textbf{\ur{ابھی}} — \emph{abhī} — \textquotedblleft{}right now\textquotedblright{}.

\begin{cousinweb}
\textbf{\ur{ابھی}} (\emph{abhī}) is \emph{ab}, \textquotedblleft{}now\textquotedblright{}, with the emphatic \emph{-hī}. CDIAL lists \emph{ab} and \emph{abhī} under Sanskrit \emph{evam eva}, \textquotedblleft{}exactly so\textquotedblright{}, crossed with a pointing \emph{a-} (CDIAL 2528).
\end{cousinweb}

\begin{grammarlens}[title={What you've built}]
One more word for this chapter.
\end{grammarlens}

\subsection*{Your turn}
\begin{itemize}
\raggedright
  \item \emph{Say it:} \emph{abhī}
  \item \emph{Say it:} \emph{abhī}, once more
  \item \emph{Say it:} say \emph{abhī}
  \item \emph{Recall:} say the Urdu for today, then the Urdu for yesterday; tomorrow, then say \emph{abhī} again
\end{itemize}

\begin{tcolorbox}[breakable,colback=teal!4,colframe=teal!35!black,title={Before you move on}]
What does \ur{ابھی} mean? (\textquotedblleft{}Right now\textquotedblright{}.) Say it once more.
\end{tcolorbox}
\section[\texorpdfstring{\ur{پھر} (phir)}{پھر (phir)}]{\ur{پھر} (phir) — then, again}

\label{lesson:UR-C76-phir}



\begin{quote}
Before the new one: say the Urdu for yesterday; tomorrow, then the Urdu for right now.
\end{quote}

\subsection*{You'll want to know: \ur{پھر}}
\textbf{\ur{پھر}} — \emph{phir} — \textquotedblleft{}then, again\textquotedblright{}.

\begin{cousinweb}
\textbf{\ur{پھر}} (\emph{phir}) is inherited: it is the \textquotedblleft{}having turned\textquotedblright{} form of a reconstructed \emph{*phirati}, \textquotedblleft{}turns\textquotedblright{}, used as an adverb (CDIAL 9078).
\end{cousinweb}

\begin{grammarlens}[title={What you've built}]
One more word for this chapter.
\end{grammarlens}

\subsection*{Your turn}
\begin{itemize}
\raggedright
  \item \emph{Say it:} \emph{phir}
  \item \emph{Say it:} \emph{phir}, once more
  \item \emph{Say it:} say \emph{phir}
  \item \emph{Recall:} say the Urdu for yesterday; tomorrow, then the Urdu for right now, then say \emph{phir} again
\end{itemize}

\begin{tcolorbox}[breakable,colback=teal!4,colframe=teal!35!black,title={Before you move on}]
What does \ur{پھر} mean? (\textquotedblleft{}Then, again\textquotedblright{}.) Say it once more.
\end{tcolorbox}
